% !TEX program = lualatex
%!TEX TS-program = lualatex
\documentclass[12pt,a4paper]{article}
\usepackage{fontspec}
\usepackage{polyglossia}
\setdefaultlanguage{polish}
\usepackage{geometry}
\geometry{margin=2.1cm}
\usepackage[table]{xcolor}
\usepackage{longtable}
\usepackage{array}
\usepackage{booktabs}
\usepackage{enumitem}
\usepackage{ragged2e}
\definecolor{tableHeader}{HTML}{E8EEF7}
\definecolor{tableStripe}{HTML}{F7F9FC}
\newcolumntype{L}[1]{>{\RaggedRight\arraybackslash}p{#1}}
\newcolumntype{C}[1]{>{\Centering\arraybackslash}p{#1}}
\renewcommand{\arraystretch}{1.22}
\setlength{\tabcolsep}{4pt}
\setlist[itemize]{leftmargin=*, itemsep=2pt, topsep=2pt}
\setlength{\parindent}{0pt}
\setlength{\parskip}{6pt}
\emergencystretch=3em
\sloppy
\begin{document}

\section*{Raport ze sprintu nr 3 - Calisia Banking App}
Data raportu: 13.05.2026

Okres sprintu: 08.05 - 13.05.2026

Product Owner / Scrum Master: Jeremiasz Kubryn

\subsection*{Lista osób}
\begin{itemize}
\item Kubryn Jeremiasz
\item Łaszuk Jakub
\item Makarevich Ksenia
\item Nuckowski Michał
\item Ruczyńska Nikola
\item Wołczko Jakub
\item Olszewski Dawid
\end{itemize}

\subsection*{Product Owner i Scrum Master sprintu}
Jeremiasz Kubryn

\subsection*{Lista historyjek na sprint}
\begingroup
\footnotesize
\rowcolors{2}{tableStripe}{white}
\begin{longtable}{@{}C{0.08\textwidth}L{0.62\textwidth}C{0.10\textwidth}L{0.12\textwidth}@{}}
\toprule
\rowcolor{tableHeader}
\textbf{Kod} & \textbf{Historyjka} & \textbf{SP} & \textbf{Status} \\
\midrule
\endfirsthead
\toprule
\rowcolor{tableHeader}
\textbf{Kod} & \textbf{Historyjka} & \textbf{SP} & \textbf{Status} \\
\midrule
\endhead
S3-01 & Logowanie e-mailem i hasłem & 8 & Done \\
S3-02 & Generowanie tokena JWT z identyfikatorem klienta i czasem wygaśnięcia & 8 & Done \\
S3-03 & Zabezpieczenie endpointów rachunków, przelewów i dashboardu atrybutem autoryzacji & 5 & Done \\
S3-04 & Zapisywanie sesji klienta w stanie aplikacji frontendu & 5 & Done \\
S3-05 & Komunikat błędu przy nieudanym logowaniu & 3 & Done \\
S3-06 & Testy logowania, generowania tokenu i odmowy dostępu & 8 & Done \\
\bottomrule
\end{longtable}
\endgroup

\subsection*{Podsumowanie sprintu}
Cel sprintu: Logowanie, JWT i ochrona API. Zapewnienie mechanizmu uwierzytelniania oraz zabezpieczenie operacji klienta przed dostępem anonimowym.

Zakres sprintu obejmował 6 historyjek o łącznej estymacji 37 SP. Wszystkie historyjki zostały dowiezione ze statusem Done, dlatego osiągnięta prędkość sprintu wyniosła 37 SP.

Udostępniono logowanie przez API, generowanie tokena JWT, dostęp do chronionych endpointów wyłącznie z poprawnym tokenem oraz obsługę komunikatów błędów przy nieudanym logowaniu. Frontend utrzymuje sesję klienta w stanie aplikacji.

Sprint znacząco podniósł bezpieczeństwo aplikacji. Od tego momentu funkcje finansowe mogą być rozwijane z założeniem, że API rozpoznaje klienta i chroni dane przed dostępem anonimowym.

\subsubsection*{Najważniejsze rezultaty}
\begin{itemize}
\item Dodano logowanie klienta z użyciem e-maila i hasła.
\item Wdrożono generowanie i walidację tokenów JWT.
\item Zabezpieczono endpointy wymagające autoryzacji.
\item Przetestowano przypadki dostępu bez tokenu i z niepoprawnymi danymi.
\end{itemize}

\subsection*{Retrospekcja sprintu}
Retrospekcja pokazała, że prace nad bezpieczeństwem wymagają dobrego przygotowania konfiguracji i testów środowiskowych.

\subsubsection*{Co się udało}
\begin{itemize}
\item Sprawna implementacja mechanizmu logowania oraz tokenów JWT.
\item Konsekwentne stosowanie autoryzacji na endpointach.
\item Efektywna współpraca backendu z frontendem przy integracji sesji.
\item Szybkie rozwiązanie problemów konfiguracyjnych JWT.
\end{itemize}
\subsubsection*{Poziom energii zespołu}
Średni poziom energii zespołu: 4,3/5.

Zespół utrzymał wysokie zaangażowanie dzięki pracy nad bezpieczeństwem aplikacji i wdrożeniem JWT.

\subsubsection*{Czego udało się nauczyć}
\begin{itemize}
\item Zespół zdobył praktyczne doświadczenie z JWT.
\item Bezpieczeństwo endpointów wymaga testów pozytywnych i negatywnych.
\item Szablony konfiguracji przyspieszają przygotowanie środowiska developerskiego.
\end{itemize}
\subsubsection*{Wnioski i działania na kolejny sprint}
\begin{itemize}
\item Przygotować szablon konfiguracji JWT w repozytorium.
\item Dopisać checklistę autoryzacji dla nowych endpointów.
\item Utrzymywać testy 401/403 przy kolejnych funkcjach klienta.
\end{itemize}

\subsection*{Następny sprint}
Product Owner / Scrum Master w następnym sprincie: Ksenia Makarevich

Główny cel następnego sprintu: Rachunki bankowe i operacje pieniężne. Stworzenie logiki kont oraz bazowych mechanizmów finansowych.

\end{document}
